% Einheit 3 von 4 – Lage als Anhang (bank/geraden/e3.jsonl, zusammenbau v0.1)
%% TODO Zeitmarke Einheit 3: Marken-Zeile nicht lesbar
%% TODO Zweigzeile Teil 1: was hier gelernt wird (Satz in Schülersprache aus dem Einheitstitel) – Einheit 3
\einheitenkopf[e3]{Einheit 3 von 4 · Lage als Anhang}
\zweigzeile{Abitur GK}
\setcounter{aufgabe}{15}

%% TODO Titel: Ich-kann-Satz fehlt in der Bank (Platzhalter: Kettenname) – Nr. 16
%% TODO Anweisung: ein Satz über der ersten Teilaufgabe fehlt in der Bank – Nr. 16
\begin{aufgabe}{Lage als Anhang}
%% TODO \rechenplatz{n}: Schrittzahl des Grundfalls fehlt in der Bank (nur bei fester Schreibform, 2.3 b)
\begin{teile}
\teil Die Richtungsvektoren der Geraden $g$ und $h$ sind Vielfache voneinander. Was ist als Nächstes zu prüfen? \kreuz{ob ein Stützpunkt von $g$ auf $h$ liegt} \\ \kreuz{ob sich die Geraden schneiden} \\ \kreuz{nichts – die Geraden sind identisch}
\teil Gegeben sind $g\colon \vec x = (2 | -1 | 3) + r \cdot (1 | 3 | 0)$ und $h\colon \vec x = (2 | -1 | 3) + t \cdot (3 | 1 | 0)$. Begründe, dass $g$ und $h$ nicht identisch sind.
\teil Die Gerade $g$ durch $P(1 | 0 | 2)$ und die Gerade $h$ durch $Q(5 | -2 | 7)$ haben beide den Richtungsvektor $(2 | -1 | 3)$. Weise nach, dass $g$ und $h$ nicht identisch sind. \hfill (Abitur 2022 GK)
\teil Ein Tunnel verläuft geradlinig durch $A(0 | 30 | 5)$ und $B(30 | 50 | 8)$, ein zweiter durch $C(10 | 0 | 2)$ und $D(70 | 40 | 8)$ (1 LE = 10 m). Prüfe, ob die beiden Tunnelachsen parallel sind.
\teil Ein Würfel hat die Ecken $A(5 | 0 | 0)$, $B(5 | 5 | 0)$, $C(0 | 5 | 0)$ und $G(0 | 5 | 5)$. Gib eine Gleichung der Geraden $h$ durch $B$ und $G$ an und begründe, dass $h$ windschief zur Geraden $AC$ ist. \hfill (Abitur 2021 GK)
\teil Die Punkte $A(1 | 1 | 0)$, $B(3 | 3 | 4)$, $C(0 | 0 | 6)$ und $D(2 | 2 | 4)$ liegen in der Ebene $E\colon x_1 - x_2 = 0$. Die Gerade $g$ verläuft durch $A$ und $B$, die Gerade $h$ durch $C$ und $D$. Begründe, ohne den Schnittpunkt zu berechnen, dass $g$ und $h$ sich schneiden müssen. \hfill (Abitur 2018 GK)
\teil Gegeben sind $g\colon \vec x = (1 | 3 | 1) + r \cdot (2 | 0 | 1)$ und die dazu echt parallele Gerade $p\colon \vec x = (1 | 5 | 1) + r \cdot (2 | 0 | 1)$. Entscheide, ob jede Gerade, die $g$ senkrecht schneidet, auch $p$ schneidet, und begründe. \hfill (Abitur 2019 GK)
\end{teile}
\end{aufgabe}

%% TODO Titel: Ich-kann-Satz fehlt in der Bank (Platzhalter: Kettenname) – Nr. 17
%% TODO Anweisung: ein Satz über der ersten Teilaufgabe fehlt in der Bank – Nr. 17
\begin{aufgabe}{Lage als Anhang – Fehler finden, Begründen}
\begin{teile}
\teil Finde den Fehler. Gegeben sind $g\colon \vec x = (4 | -1 | 2) + r \cdot (1 | 2 | -1)$ und $h\colon \vec x = (4 | -1 | 2) + t \cdot (2 | 1 | 1)$. Ben sagt: „Beide Geraden haben denselben Stützpunkt, also sind sie identisch.“
\teil Begründe, warum zwei Geraden, die in einer gemeinsamen Ebene liegen und nicht parallel sind, sich schneiden müssen.
\end{teile}
\end{aufgabe}
